\section{Ellipse d'action, radian et double retour d'holonomie}
\label{sec:elipse-accion-holonomia-rev8}
\index{action!ellipse}
\index{radian!secteur d'action}
\index{Möbius!double retour}

Le théorème focal précédent fixe la forme angulaire et le lecteur déterminantiel
fixe l'échelle d'action. Définissons
\[
 \eta=\operatorname{artanh}\frac{C^*}{A},
 \qquad
 a_S=\sqrt{2\hbar_{\rm HMT}^{\rm ret}e^\eta},
 \qquad
 b_S=\sqrt{2\hbar_{\rm HMT}^{\rm ret}e^{-\eta}}.
\]
L'orbite symplectique
\[
 Q(\theta)=a_S\cos\theta,
 \qquad
 P(\theta)=b_S\sin\theta
\]
conserve simultanément le rapport de forme du bloc angulaire et l'aire
fixée par la droite d'action.

\begin{theorem}[Aire balayée par phase]
\label{thm:area-fase-elipse-rev8}
Pour tout \(\theta\),
\[
 \operatorname{Area}(\theta)
 =\frac12\int_0^\theta(QP'-PQ')\,dt
 =\hbar_{\rm HMT}^{\rm ret}\theta.
\]
Par conséquent,
\[
 \operatorname{Area}(1\ {\rm rad})=\hbar_{\rm HMT}^{\rm ret},
 \qquad
 \operatorname{Area}(2\pi)=h_{\rm HMT}^{\rm ret}.
\]
\end{theorem}

\begin{proof}
Comme \(QP'-PQ'=a_Sb_S\) et
\(a_Sb_S=2\hbar_{\rm HMT}^{\rm ret}\), l'intégrale vaut
\(\hbar_{\rm HMT}^{\rm ret}\theta\). La seconde identité utilise
\(h_{\rm HMT}^{\rm ret}=2\pi\hbar_{\rm HMT}^{\rm ret}\).
\end{proof}

APP fournit l'aire orientée ; \(A,C^*\) fixent la forme et les foyers ; la branche
déterminantielle fixe l'échelle. Le tour primitif enferme \(h\) et un radian
de phase balaie \(\hbar\). Le radian ordinaire reste
\(1\,\mathrm{rad}=180^\circ/\pi\) ; la construction ajoute sa généalogie comme
secteur unitaire d'action, non une autre constante angulaire.

Pour une section persistante d'holonomie
\(\varepsilon_\gamma\in\{+1,-1\}\), la condition
\[
 \varepsilon_\gamma e^{iJ_\gamma/\hbar}=1
\]
distingue deux clôtures. Dans le cylindre, \(\varepsilon_\gamma=+1\) et le cycle
primitif transporte \(J=2\pi\hbar=h\). Dans le ruban de Möbius,
\(\varepsilon_\gamma=-1\) : un tour visible change de feuillet et transporte
\(J=\pi\hbar=h/2\) ; le second rétablit l'orientation et complète \(h\). Par
radian de la phase interne, on transporte \(\hbar\) dans les deux cas ; le facteur
\(1/2\) appartient à la projection sur le tour visible de la base de Möbius.

\begin{figure}[htbp]
\centering
\resizebox{.98\textwidth}{!}{\begin{tikzpicture}[
  font=\scriptsize,
  box/.style={draw,rounded corners=2mm,align=center,minimum height=.9cm,
    inner xsep=5pt},
  arr/.style={-{Stealth[length=2.2mm]},thick}]

  \begin{scope}[shift={(0,0)}]
    \draw[->,hmtgray] (-3.0,0)--(3.0,0) node[right]{\(u_+\)};
    \draw[->,hmtgray] (0,-2.0)--(0,2.0) node[above]{\(u_-\)};
    \draw[hmtblue,very thick] (0,0) ellipse (2.65cm and 1.45cm);
    \fill[hmtgold] (-1.85,0) circle (2pt) node[below]{\(F_-^{(\lambda)}\)};
    \fill[hmtgold] (1.85,0) circle (2pt) node[below]{\(F_+^{(\lambda)}\)};
    \node[text=hmtblue,font=\bfseries] at (0,2.55) {ellipse des taux};
    \node[align=center,text=hmtgray] at (0,-2.4)
      {\(q_+=e^{-a_\lambda^2}\), \(q_-=e^{-b_\lambda^2}\)};
  \end{scope}

  \draw[arr,hmtgold] (3.3,0)--node[above,align=center]{même forme\\échelle d’action} (6.0,0);

  \begin{scope}[shift={(9.0,0)}]
    \draw[->,hmtgray] (-2.8,0)--(2.8,0) node[right]{\(Q\)};
    \draw[->,hmtgray] (0,-2.0)--(0,2.0) node[above]{\(P\)};
    \fill[hmtlightgold] (0,0)--(2.5,0)
      arc[start angle=0,end angle=57.2958,x radius=2.5cm,y radius=1.35cm]--cycle;
    \draw[hmtviolet,very thick] (0,0) ellipse (2.5cm and 1.35cm);
    \draw[hmtgold,thick] (0,0)--(2.5,0);
    \draw[hmtgold,thick] (0,0)--({2.5*cos(57.2958)},{1.35*sin(57.2958)});
    \node[text=hmtviolet,font=\bfseries] at (0,2.55) {ellipse d’action};
    \node[align=center] at (.7,.48) {\(\theta=1\)\\aire \(\hbar\)};
    \node[align=center,text=hmtgray] at (0,-2.4) {tour \(2\pi\) : aire \(h\)};
  \end{scope}

  \node[box,draw=hmtblue,fill=hmtlightblue,minimum width=3.2cm] (cyl) at (3.2,-4.3)
    {cilindro \(\varepsilon=+1\)\\\(2\pi\mapsto h\)};
  \node[box,draw=hmtviolet,fill=hmtlightviolet,minimum width=4.2cm] (mob) at (9.0,-4.3)
    {Möbius \(\varepsilon=-1\)\\\(2\pi\mapsto h/2\), \(4\pi\mapsto h\)};
  \draw[arr,hmtgray] (cyl)--node[
    above,align=center,text width=1.9cm,fill=white,inner sep=1pt
  ]{changement d’\\holonomie} (mob);
  \node[align=center,text=hmtgray,font=\footnotesize] at (6.1,-5.55)
    {La phase interne transporte \(\hbar\) par radian ;\\le facteur deux appartient au tour visible.};
\end{tikzpicture}
}
\caption{L'ellipse des taux fixe la forme ; son relèvement à la droite d'action fixe
l'aire. La clôture cylindrique et le double retour de Möbius distinguent la phase
interne du tour visible.}
\label{fig:elipse-accion-holonomia-rev8}
\end{figure}

\subsection{Géométrie de Hilbert--Beltrami--Klein dans l'ellipse d'action}
\label{sec:hilbert-beltrami-klein-elipse-accion}

La même ellipse qui porte l'aire symplectique admet une géométrie projective
intrinsèque. Soit
\[
 \mathbb E_{\rm act}
 :=\left\{(Q,P)\in\mathbb R^2:
 \frac{Q^2}{a_S^2}+\frac{P^2}{b_S^2}<1\right\},
 \qquad
 \mathcal L_{\rm act}(Q,P)
 :=\left(\frac{Q}{a_S},\frac{P}{b_S}\right).
\]
L'application linéaire \(\mathcal L_{\rm act}\) identifie
\(\mathbb E_{\rm act}\) au disque unité \(\mathbb D\). Pour deux points
distincts \(x,y\in\mathbb E_{\rm act}\), soient \(a,b\in\partial
\mathbb E_{\rm act}\) les intersections de la droite \(xy\) avec la frontière,
ordonnées comme \(a,x,y,b\). Nous définissons
\begin{equation}
 d_{\rm H}^{\rm act}(x,y)
 :=\frac12\log
 \frac{\lVert y-a\rVert\,\lVert x-b\rVert}
      {\lVert x-a\rVert\,\lVert y-b\rVert}.
 \label{eq:distancia-hilbert-elipse-accion}
\end{equation}
L'expression est le birapport des quatre points colinéaires et est donc
indépendante de la norme auxiliaire utilisée sur cette droite.

\begin{proposition}[Réalisation de Beltrami--Klein de l'ellipse d'action]
\label{prop:beltrami-klein-elipse-accion}
Le couple \((\mathbb E_{\rm act},d_{\rm H}^{\rm act})\) est isométrique, par
\(\mathcal L_{\rm act}\), au modèle de Beltrami--Klein du plan hyperbolique.
Ses géodésiques non orientées sont exactement les cordes ouvertes de
l'ellipse. La frontière \(\partial\mathbb E_{\rm act}\) conserve, de manière
simultanée et typée, l'orbite symplectique
\[
 \gamma_{\rm act}(\theta)
 =(a_S\cos\theta,b_S\sin\theta),
 \qquad
 \frac12\int_0^\theta(QP'-PQ')\,dt
 =\hbar_{\rm HMT}^{\rm ret}\theta.
\]
\end{proposition}

\begin{proof}
Les transformations projectives préservent les birapports. Comme
\(\mathcal L_{\rm act}\) est une transformation linéaire inversible qui envoie
l'ellipse sur le disque unité, elle transporte
\eqref{eq:distancia-hilbert-elipse-accion} vers la distance de Hilbert du
disque. Celle-ci est la métrique de Beltrami--Klein, dont les géodésiques sont les cordes.
L'identité symplectique de la frontière est le
\cref{thm:area-fase-elipse-rev8} écrit dans la carte normalisée.
\end{proof}

La carte complète est donc typée par
\[
 \mathfrak R_{\rm act}:
 (A,C^*,\hbar_{\rm HMT}^{\rm ret})
 \longmapsto
 (\mathbb E_{\rm act},d_{\rm H}^{\rm act},
  dP\wedge dQ,\gamma_{\rm act}).
\]
La métrique projective gouverne l'intérieur ; la forme symplectique et la loi
d'aire gouvernent la frontière parcourue. Une réalisation comme région physique ou
comme domaine de matrices de covariance se formule par un foncteur ultérieur
qui conserve explicitement ces deux structures ; la construction précédente
fixe déjà son domaine mathématique et les invariants que ce foncteur doit préserver.
